\section{第 6 章 系统测试与结果分析}

\subsection{6.1 测试目标与方法}

本章测试工作的重点，不只是说明系统页面能够打开或接口能够返回结果，更重要的是验证前文提出的多条业务链路是否真正形成闭环，并能够输出具有参考意义的分析结果。结合当前项目实际情况，系统测试主要围绕四个方面展开：一是主要页面与接口能否稳定访问；二是关键分析逻辑是否符合设计预期；三是机器学习误差校正结果是否真实、可解释；四是在外部数据波动或参数异常时，系统是否具备合理回退机制。

系统测试方式为自动化测试加页面级人工检查。自动化部分主要使用 pytest 与 Flask test\_client 对历史气候、探空分析、预报时次识别等关键逻辑进行验证；页面层面则对 /advanced-forecast、/upperair、/history/yearly、/agro-dashboard 等主要模块进行了功能检查和结果核对。对于机器学习模块，还保留了 training\_manifest.json 和 metrics\_report.json 两类评估产物，用于支撑模型训练过程与测试结果的追溯。

\subsection{6.2 主要功能流程验证}

从整体流程上看，系统首页可以作为统一的导航入口，把用户引导到多模式预报、探空分析、历史气候、农业气象等主要模块上。各个模块虽然业务侧重点不同，但是都以同一个杭州站点口径、统一的页面导航结构以及相对稳定的的时间组织规则为基础，从而使得整个系统具有较好的整体性。

多模式预报模块测试的重点是72小时精细化预报、多模式对比和ECMWF机器学习校正链路能否协同工作；历史气候模块重点检验原始极值和EW指数模式切换、四季代表事件展示、页面渲染稳定度；探空模块重点检验输入时间转换、数据抓取、参数解析、异常值表达的一致性；农业模块重点检验总览页、作物专题页、AI建议接口之间能否形成作物状态、风险提示、农事建议的闭环。整体测试结果说明各个主要功能模块都可以按照设计意图实现协同工作，系统主业务流程是连续的。

\subsection{6.3 关键逻辑验证结果}

历史气候模块的验证主要看EW指数代表事件筛选逻辑。自动化测试结果表明，系统可以正确地从历史同月同日的样本中提取出来，对低温事件采用反向百分位处理，对于同一个低温、高温或者降水过程中的连续异常日进行合并，防止同一个过程出现多次上榜的情况。同时EW的结果可以按照春、夏、秋、冬四季分组显示，并且控制每个季节每个要素的代表性事件的数量，说明本模块有一定的统计解释力，而不是只停留在图表的展示上。

探空模块主要的验证就是时次转换以及异常参数的处理。测试结果表明，页面输入的北京时间能够正确换算为 UTC 标准探空时次；当用户输入 2026-04-10 08:00 时，系统可正确解析为 2026-04-10 00:00 UTC。接口元数据中可以返回请求本地时次、请求UTC时次、实际解析时次和是否使用回退。当CAPE出现负值或者计算出错的时候，系统会统一用N/A来处理，不会直接显示负数，从而降低了给用户造成误导的风险。

从预报时次统一下手看，经过相关测试发现，新版HZ\_forecast\_*.csv文件的时次识别规则可以保持稳定工作，凌晨04:15左右采集的样本会被识别成前一天20:00本地起报，对应的周期是12Z；下午16:15左右采集的样本会被识别成当天08:00本地起报，对应的周期是00Z；明显偏离支持窗口的异常样本会被剔除。这就意味着训练样本的组织方式和页面运行时的口径是相同的，没有造成隐性的偏差。

\subsection{6.4 模型评估结果分析}

表 6-1  机器学习误差校正连续变量评估结果

表 6-2  机器学习误差校正降水评估结果

机器学习误差校正模块最终的评价结果说明，更一致的站点口径、更严格的测试划分条件下，连续变量有比较明显的改善。温度MAE由原来的2.355°C降到现在的1.832°C，湿度MAE由原来的9.851\%降到现在的9.085\%，风速MAE由原来的4.413m/s降到现在的2.613m/s。从分时效结果可以看出，短时效内改善的更加明显，后处理模型对短中期预报误差有较好的修正作用。

降水变量的结果就更复杂了。分类指标并没有比原来的预报更好，但是Recall、F1基本不变，Precision基本不变或略有上升，雨样本MAE和全样本MAE都明显降低。由此可以看出目前降水模型策略偏向保守，即减少误报降水，但是对于成功识别出的降水，雨量级估计更加接近实际情况。对农业灌溉、水库调度、排涝准备等场景来说，量级误差的改善仍有较高的应用价值。

为了更直观地比较不同的版本误差校正方案在严格时间隔离测试集上的表现，本文在原始预报、V1 和 V2 三种结果之间又增加了统一对比，见表 6-3。需要说明的是，该表对应最后的测试集使用 issue\_time 留出，测试起报时间为 2026年4月4日到2026年4月10日。

从结果可知，本轮V2实验并没有比V1更好，这是合理的。在目前的数据规模大约为 2 个月的时候，随机森林这种 bagging 模型一般比 boosting 模型要更加稳定，对局部噪声以及样本不平衡性都比较敏感。其次Boosting模型一般需要更多的样本、更多的特征工程以及更加细致的参数调节才能发挥出它的优势，目前的研究阶段主要是集中于站点口径统一、issue\_time规则统一以及严格的时段隔离评估。因此本文的核心结论并不是V2一定比V1好，而是说明在小样本量的情况下，随机森林仍然是误差校正的最佳选择；而和只换模型相比，严格控制时间切分、防止目标时刻泄漏、保证训练和运行时口径一致，往往对结果的可信度影响更大。

表 6-3  不同版本误差校正效果对比（严格时间隔离测试集）

注：加粗为最优结果。V1 在温度、湿度、风速、降水量上均优于 Raw 和 V2。

\subsection{6.5 结果讨论}

综合本章测试可以看出，本系统已经形成较完整的业务闭环。多模式预报、探空分析、历史气候、机器学习校正和农业气象服务并非孤立模块，而是在统一站点口径、统一页面入口和统一时间规则下被组织到同一平台之中。测试结果说明，这种组织方式在当前实现中是可运行且基本稳定的。

另一方面，系统测试也表明，本项目真正难的是页面本身，而不是时间规则、统计逻辑和异常处理等容易被忽略的细节。issue\_time 推断、北京时间到UTC的转换、EW榜单去重、CAPE异常值处理等都会影响结果的可信度。目前这些重要细节都已经得到有针对性的验证，所以本系统的结果比以前更加真实、有参考价值。

